\chapter*{第四部分学习路线}
\addcontentsline{toc}{chapter}{第四部分学习路线：把物理图像扩展到光、热、原子与实验}
\markboth{第四部分学习路线}{第四部分学习路线}

\begin{learningbox}{把物理图像扩展到光、热、原子与实验}
第四部分把前面建立的模型意识扩展到光、热、原子与实验：光学看光路，热学看微观模型和状态图，近代物理看实验事实，实验专题训练测量、误差和图像处理。
\end{learningbox}

\section*{一、这一部分的阅读顺序}

\begin{enumerate}[itemsep=0.2em]
  \item \textbf{Day 25 光学：} 先画光路，再用反射、折射、成像和干涉规律。
  \item \textbf{Day 26 热学：} 先分清状态量和过程量，再看气体状态图。
  \item \textbf{Day 27 近代物理：} 先记关键实验事实，再理解经典图像为什么要修正。
  \item \textbf{Day 28 实验：} 先看目的和原理公式，再看测量量、斜率和误差。
  \item \textbf{Day 29 到 Day 30：} 回到全书模型，训练综合判断。
\end{enumerate}

\section*{二、这一部分最重要的能力}

\begin{itemize}[itemsep=0.2em]
  \item \textbf{把图画准确：} 光路图、$p$-$V$ 图、能级图、实验电路图都直接决定解题方向。
  \item \textbf{尊重实验事实：} 光电效应、卢瑟福散射、原子能级不能只靠生活直觉推断。
  \item \textbf{会读坐标图：} 斜率、面积、截距和变化趋势常常比单个数据更重要。
  \item \textbf{会复盘模型：} 综合题是在合适位置切换熟悉模型。
\end{itemize}

\section*{三、卡住时先回到哪里}

光学题卡住，先重画法线和光路；热学题卡住，先判断是状态变化还是过程传热、做功；近代物理题卡住，先回到实验事实和能量守恒；实验和综合题卡住，先问测什么、画什么图。
